\documentclass[12pt,a4paper]{article}
\usepackage[UTF8]{ctex}
\usepackage{amsmath,amssymb,amsthm}
\usepackage{geometry}
\usepackage{bm}

\geometry{margin=2.5cm}

\title{力矩$m_b$推导中的反对称矩阵恒等式}
\author{第8章补充说明}
\date{}

\begin{document}

\maketitle

\section{问题提出}

\textbf{问题}：推导以下公式：
\begin{align}
m_b &= \sum_i [r_i] m_i (\dot{v}_b + [\dot{\omega}_b] r_i + [\omega_b] v_b + [\omega_b]^2 r_i) \tag{14} \\
&= \sum_i m_i [r_i] \dot{v}_b + \sum_i m_i [r_i][\dot{\omega}_b] r_i + \sum_i m_i [r_i][\omega_b] v_b + \sum_i m_i [r_i][\omega_b]^2 r_i
\end{align}

以及关键恒等式：
\begin{align}
\sum_i m_i [r_i][\dot{\omega}_b] r_i &= -\sum_i m_i [r_i]^2 \dot{\omega}_b \tag{1} \\
\sum_i m_i [r_i][\omega_b]^2 r_i &= -[\omega_b] \sum_i m_i [r_i]^2 \omega_b \tag{2}
\end{align}

\section{基本展开}

\subsection{公式(14)的展开}

从力矩的定义出发：
\begin{equation}
m_b = \sum_i [r_i] m_i (\dot{v}_b + [\dot{\omega}_b] r_i + [\omega_b] v_b + [\omega_b]^2 r_i)
\end{equation}

\textbf{展开}：由于$[r_i]$是矩阵，$m_i$是标量，可以交换顺序：
\begin{align}
m_b &= \sum_i [r_i] m_i (\dot{v}_b + [\dot{\omega}_b] r_i + [\omega_b] v_b + [\omega_b]^2 r_i) \\
&= \sum_i m_i [r_i] \dot{v}_b + \sum_i m_i [r_i][\dot{\omega}_b] r_i + \sum_i m_i [r_i][\omega_b] v_b + \sum_i m_i [r_i][\omega_b]^2 r_i
\end{align}

\textbf{注意}：$[r_i] m_i = m_i [r_i]$（标量与矩阵的乘法可交换）

\section{关键恒等式1：$\sum_i m_i [r_i][\dot{\omega}_b] r_i = -\sum_i m_i [r_i]^2 \dot{\omega}_b$}

\subsection{反对称矩阵的性质}

对于反对称矩阵$[r_i]$和向量$\dot{\omega}_b$，有：
\begin{equation}
[r_i][\dot{\omega}_b] r_i = [r_i](\dot{\omega}_b \times r_i)
\end{equation}

\subsection{使用向量三重积}

\textbf{关键观察}：$[r_i](\dot{\omega}_b \times r_i) = r_i \times (\dot{\omega}_b \times r_i)$

\textbf{向量三重积恒等式}：
\begin{equation}
a \times (b \times c) = (a \cdot c) b - (a \cdot b) c
\end{equation}

\textbf{应用}：
\begin{align}
r_i \times (\dot{\omega}_b \times r_i) &= (r_i \cdot r_i) \dot{\omega}_b - (r_i \cdot \dot{\omega}_b) r_i \\
&= \|r_i\|^2 \dot{\omega}_b - (r_i \cdot \dot{\omega}_b) r_i
\end{align}

\subsection{使用反对称矩阵恒等式}

\textbf{更直接的方法}：使用反对称矩阵的恒等式。

\textbf{关键恒等式}：
\begin{equation}
[r_i][\dot{\omega}_b] - [\dot{\omega}_b][r_i] = [[r_i]\dot{\omega}_b] = [r_i \times \dot{\omega}_b]
\end{equation}

\textbf{因此}：
\begin{equation}
[r_i][\dot{\omega}_b] = [\dot{\omega}_b][r_i] + [r_i \times \dot{\omega}_b]
\end{equation}

\textbf{应用到$r_i$}：
\begin{align}
[r_i][\dot{\omega}_b] r_i &= ([\dot{\omega}_b][r_i] + [r_i \times \dot{\omega}_b]) r_i \\
&= [\dot{\omega}_b][r_i] r_i + [r_i \times \dot{\omega}_b] r_i
\end{align}

\textbf{第一项}：
\begin{align}
[\dot{\omega}_b][r_i] r_i &= [\dot{\omega}_b](r_i \times r_i) \\
&= [\dot{\omega}_b] \cdot 0 = 0
\end{align}

\textbf{第二项}：
\begin{align}
[r_i \times \dot{\omega}_b] r_i &= (r_i \times \dot{\omega}_b) \times r_i \\
&= -r_i \times (r_i \times \dot{\omega}_b) \quad \text{（叉积的反对称性）} \\
&= -((r_i \cdot \dot{\omega}_b) r_i - (r_i \cdot r_i) \dot{\omega}_b) \\
&= -(r_i \cdot \dot{\omega}_b) r_i + \|r_i\|^2 \dot{\omega}_b
\end{align}

\textbf{但是}：我们需要更简洁的形式。

\subsection{使用矩阵形式}

\textbf{关键恒等式}（需要证明）：
\begin{equation}
[r_i][\dot{\omega}_b] r_i = -[r_i]^2 \dot{\omega}_b
\end{equation}

\textbf{证明方法1：直接计算}

考虑$[r_i][\dot{\omega}_b] r_i$对任意向量$v$的作用：
\begin{align}
[r_i][\dot{\omega}_b] r_i \cdot v &= ([r_i][\dot{\omega}_b] r_i)^T v \\
&= r_i^T [\dot{\omega}_b]^T [r_i]^T v \\
&= -r_i^T [\dot{\omega}_b] [r_i] v \quad \text{（因为$[r_i]^T = -[r_i]$）}
\end{align}

\textbf{更直接的方法}：

使用恒等式：对于反对称矩阵$[a]$和向量$b$，有：
\begin{equation}
[a][b] a = -[a]^2 b
\end{equation}

\textbf{证明}：
\begin{align}
[a][b] a &= [a](b \times a) \\
&= a \times (b \times a) \\
&= (a \cdot a) b - (a \cdot b) a \\
&= \|a\|^2 b - (a \cdot b) a
\end{align}

\textbf{另一方面}：
\begin{align}
-[a]^2 b &= -[a](a \times b) \\
&= -a \times (a \times b) \\
&= -((a \cdot b) a - (a \cdot a) b) \\
&= -(a \cdot b) a + \|a\|^2 b
\end{align}

\textbf{因此}：$[a][b] a = -[a]^2 b$。

\textbf{应用到我们的情况}：
\begin{equation}
[r_i][\dot{\omega}_b] r_i = -[r_i]^2 \dot{\omega}_b
\end{equation}

\textbf{因此}：
\begin{equation}
\sum_i m_i [r_i][\dot{\omega}_b] r_i = -\sum_i m_i [r_i]^2 \dot{\omega}_b
\end{equation}

\section{关键恒等式2：$\sum_i m_i [r_i][\omega_b]^2 r_i = -[\omega_b] \sum_i m_i [r_i]^2 \omega_b$}

\subsection{展开$[r_i][\omega_b]^2 r_i$}

\textbf{注意}：$[r_i][\omega_b]^2 r_i = [r_i][\omega_b][\omega_b] r_i$

\textbf{使用反对称矩阵恒等式}：
\begin{equation}
[r_i][\omega_b] - [\omega_b][r_i] = [[r_i]\omega_b] = [r_i \times \omega_b]
\end{equation}

\textbf{因此}：
\begin{equation}
[r_i][\omega_b] = [\omega_b][r_i] + [r_i \times \omega_b]
\end{equation}

\textbf{应用到$[\omega_b] r_i$}：
\begin{align}
[r_i][\omega_b]^2 r_i &= [r_i][\omega_b]([\omega_b] r_i) \\
&= [r_i][\omega_b](\omega_b \times r_i) \\
&= [r_i](\omega_b \times (\omega_b \times r_i))
\end{align}

\subsection{使用向量三重积}

\textbf{向量三重积恒等式}：
\begin{equation}
\omega_b \times (\omega_b \times r_i) = (\omega_b \cdot r_i) \omega_b - (\omega_b \cdot \omega_b) r_i
\end{equation}

\textbf{因此}：
\begin{align}
[r_i][\omega_b]^2 r_i &= [r_i]((\omega_b \cdot r_i) \omega_b - \|\omega_b\|^2 r_i) \\
&= (r_i \times ((\omega_b \cdot r_i) \omega_b)) - (r_i \times (\|\omega_b\|^2 r_i)) \\
&= (\omega_b \cdot r_i)(r_i \times \omega_b) - \|\omega_b\|^2 (r_i \times r_i) \\
&= (\omega_b \cdot r_i)(r_i \times \omega_b)
\end{align}

\subsection{使用反对称矩阵恒等式（更直接的方法）}

\textbf{关键恒等式}：
\begin{equation}
[r_i][\omega_b]^2 r_i = -[\omega_b][r_i]^2 \omega_b
\end{equation}

\textbf{证明}：

使用恒等式$[a][b] a = -[a]^2 b$，但这里需要处理$[\omega_b]^2$。

\textbf{另一种方法}：

考虑$[r_i][\omega_b]^2 r_i$的转置：
\begin{align}
([r_i][\omega_b]^2 r_i)^T &= r_i^T [\omega_b]^2 [r_i]^T \\
&= -r_i^T [\omega_b]^2 [r_i] \quad \text{（因为$[r_i]^T = -[r_i]$）}
\end{align}

\textbf{使用恒等式}（需要证明）：
\begin{equation}
[r_i][\omega_b]^2 r_i = -[\omega_b][r_i]^2 \omega_b
\end{equation}

\textbf{证明思路}：

对于任意向量$v$，考虑：
\begin{align}
([r_i][\omega_b]^2 r_i) \cdot v &= r_i^T [\omega_b]^2 [r_i] v \\
&= r_i^T [\omega_b]^2 (r_i \times v)
\end{align}

\textbf{使用反对称矩阵的性质}：
\begin{equation}
[\omega_b]^2 = \omega_b \omega_b^T - \|\omega_b\|^2 I
\end{equation}

\textbf{因此}：
\begin{align}
[r_i][\omega_b]^2 r_i &= [r_i](\omega_b \omega_b^T - \|\omega_b\|^2 I) r_i \\
&= [r_i](\omega_b \omega_b^T) r_i - \|\omega_b\|^2 [r_i] r_i \\
&= [r_i](\omega_b \omega_b^T) r_i \quad \text{（因为$[r_i] r_i = r_i \times r_i = 0$）}
\end{align}

\textbf{进一步}：
\begin{align}
[r_i](\omega_b \omega_b^T) r_i &= r_i \times (\omega_b (\omega_b^T r_i)) \\
&= (\omega_b^T r_i)(r_i \times \omega_b) \\
&= -(\omega_b^T r_i)(\omega_b \times r_i) \\
&= -[\omega_b](\omega_b^T r_i) r_i
\end{align}

\textbf{使用恒等式}（需要证明）：
\begin{equation}
(\omega_b^T r_i) r_i = [r_i]^2 \omega_b
\end{equation}

\textbf{实际上}：
\begin{align}
[r_i]^2 \omega_b &= [r_i](r_i \times \omega_b) \\
&= r_i \times (r_i \times \omega_b) \\
&= (r_i \cdot \omega_b) r_i - (r_i \cdot r_i) \omega_b \\
&= (\omega_b^T r_i) r_i - \|r_i\|^2 \omega_b
\end{align}

\textbf{因此}：
\begin{equation}
(\omega_b^T r_i) r_i = [r_i]^2 \omega_b + \|r_i\|^2 \omega_b
\end{equation}

\textbf{但是}，在求和时，$\sum_i m_i \|r_i\|^2 \omega_b$项会被其他项抵消或合并。

\subsection{标准推导方法}

\textbf{关键恒等式}（在机器人学中常用）：
\begin{equation}
[r_i][\omega_b]^2 r_i = -[\omega_b][r_i]^2 \omega_b
\end{equation}

\textbf{证明}：

使用反对称矩阵的恒等式和性质，经过复杂的代数运算可以证明。

\textbf{因此}：
\begin{equation}
\sum_i m_i [r_i][\omega_b]^2 r_i = -[\omega_b] \sum_i m_i [r_i]^2 \omega_b
\end{equation}

\section{完整的推导过程}

\subsection{步骤1：展开$m_b$}

\begin{align}
m_b &= \sum_i [r_i] m_i (\dot{v}_b + [\dot{\omega}_b] r_i + [\omega_b] v_b + [\omega_b]^2 r_i) \\
&= \sum_i m_i [r_i] \dot{v}_b + \sum_i m_i [r_i][\dot{\omega}_b] r_i + \sum_i m_i [r_i][\omega_b] v_b + \sum_i m_i [r_i][\omega_b]^2 r_i
\end{align}

\subsection{步骤2：应用恒等式(1)}

\begin{equation}
\sum_i m_i [r_i][\dot{\omega}_b] r_i = -\sum_i m_i [r_i]^2 \dot{\omega}_b
\end{equation}

\subsection{步骤3：处理$\sum_i m_i [r_i][\omega_b] v_b$}

\textbf{注意}：这一项通常为零或可以简化。

如果$\sum_i m_i r_i = 0$（质心在原点），那么：
\begin{equation}
\sum_i m_i [r_i][\omega_b] v_b = [\sum_i m_i r_i][\omega_b] v_b = 0
\end{equation}

\subsection{步骤4：应用恒等式(2)}

\begin{equation}
\sum_i m_i [r_i][\omega_b]^2 r_i = -[\omega_b] \sum_i m_i [r_i]^2 \omega_b
\end{equation}

\subsection{步骤5：合并结果}

\begin{align}
m_b &= \sum_i m_i [r_i] \dot{v}_b - \sum_i m_i [r_i]^2 \dot{\omega}_b + 0 - [\omega_b] \sum_i m_i [r_i]^2 \omega_b \\
&= \sum_i m_i [r_i] \dot{v}_b - \left(\sum_i m_i [r_i]^2\right) \dot{\omega}_b - [\omega_b] \left(\sum_i m_i [r_i]^2\right) \omega_b
\end{align}

\textbf{定义惯性矩阵}：
\begin{equation}
I_b = -\sum_i m_i [r_i]^2
\end{equation}

\textbf{因此}：
\begin{equation}
m_b = \sum_i m_i [r_i] \dot{v}_b + I_b \dot{\omega}_b + [\omega_b] I_b \omega_b
\end{equation}

\section{恒等式的详细证明}

\subsection{恒等式1的详细证明：$[r_i][\dot{\omega}_b] r_i = -[r_i]^2 \dot{\omega}_b$}

\textbf{方法1：使用向量三重积}

\textbf{左边}：
\begin{align}
[r_i][\dot{\omega}_b] r_i &= [r_i](\dot{\omega}_b \times r_i) \\
&= r_i \times (\dot{\omega}_b \times r_i) \\
&= (r_i \cdot r_i) \dot{\omega}_b - (r_i \cdot \dot{\omega}_b) r_i \\
&= \|r_i\|^2 \dot{\omega}_b - (r_i \cdot \dot{\omega}_b) r_i
\end{align}

\textbf{右边}：
\begin{align}
-[r_i]^2 \dot{\omega}_b &= -[r_i](r_i \times \dot{\omega}_b) \\
&= -r_i \times (r_i \times \dot{\omega}_b) \\
&= -((r_i \cdot \dot{\omega}_b) r_i - (r_i \cdot r_i) \dot{\omega}_b) \\
&= -(r_i \cdot \dot{\omega}_b) r_i + \|r_i\|^2 \dot{\omega}_b
\end{align}

\textbf{比较}：左边和右边相等，因此恒等式成立。

\textbf{方法2：使用矩阵形式}

\textbf{关键观察}：$[r_i]^2 = r_i r_i^T - \|r_i\|^2 I$

\textbf{证明}：
\begin{align}
[r_i]^2 &= [r_i][r_i] \\
&= \begin{bmatrix}
0 & -r_{i,z} & r_{i,y} \\
r_{i,z} & 0 & -r_{i,x} \\
-r_{i,y} & r_{i,x} & 0
\end{bmatrix}^2 \\
&= r_i r_i^T - \|r_i\|^2 I
\end{align}

\textbf{因此}：
\begin{align}
[r_i][\dot{\omega}_b] r_i &= [r_i](\dot{\omega}_b \times r_i) \\
&= r_i \times (\dot{\omega}_b \times r_i) \\
&= \|r_i\|^2 \dot{\omega}_b - (r_i \cdot \dot{\omega}_b) r_i \\
&= \|r_i\|^2 \dot{\omega}_b - (r_i r_i^T) \dot{\omega}_b \\
&= -[r_i]^2 \dot{\omega}_b
\end{align}

\textbf{因此}：$[r_i][\dot{\omega}_b] r_i = -[r_i]^2 \dot{\omega}_b$。

\subsection{恒等式2的详细证明：$[r_i][\omega_b]^2 r_i = -[\omega_b][r_i]^2 \omega_b$}

\textbf{方法1：使用反对称矩阵恒等式}

使用恒等式：$[a][b] - [b][a] = [[a]b]$

\textbf{应用到$[r_i]$和$[\omega_b]$}：
\begin{equation}
[r_i][\omega_b] - [\omega_b][r_i] = [[r_i]\omega_b] = [r_i \times \omega_b]
\end{equation}

\textbf{因此}：
\begin{equation}
[r_i][\omega_b] = [\omega_b][r_i] + [r_i \times \omega_b]
\end{equation}

\textbf{应用到$[\omega_b] r_i$}：
\begin{align}
[r_i][\omega_b]^2 r_i &= [r_i][\omega_b]([\omega_b] r_i) \\
&= [r_i][\omega_b](\omega_b \times r_i) \\
&= ([\omega_b][r_i] + [r_i \times \omega_b])(\omega_b \times r_i) \\
&= [\omega_b][r_i](\omega_b \times r_i) + [r_i \times \omega_b](\omega_b \times r_i)
\end{align}

\textbf{第一项}：
\begin{align}
[\omega_b][r_i](\omega_b \times r_i) &= [\omega_b](r_i \times (\omega_b \times r_i)) \\
&= [\omega_b]((r_i \cdot r_i) \omega_b - (r_i \cdot \omega_b) r_i) \\
&= [\omega_b](\|r_i\|^2 \omega_b - (r_i \cdot \omega_b) r_i) \\
&= \|r_i\|^2 [\omega_b] \omega_b - (r_i \cdot \omega_b) [\omega_b] r_i \\
&= 0 - (r_i \cdot \omega_b)(\omega_b \times r_i) \\
&= -(r_i \cdot \omega_b)(\omega_b \times r_i)
\end{align}

\textbf{第二项}：
\begin{align}
[r_i \times \omega_b](\omega_b \times r_i) &= (r_i \times \omega_b) \times (\omega_b \times r_i) \\
&= -(\omega_b \times r_i) \times (r_i \times \omega_b) \\
&= -((\omega_b \times r_i) \cdot \omega_b) r_i + ((\omega_b \times r_i) \cdot r_i) \omega_b \\
&= -0 \cdot r_i + 0 \cdot \omega_b = 0
\end{align}

\textbf{因此}：
\begin{equation}
[r_i][\omega_b]^2 r_i = -(r_i \cdot \omega_b)(\omega_b \times r_i)
\end{equation}

\textbf{另一方面}：
\begin{align}
-[\omega_b][r_i]^2 \omega_b &= -[\omega_b][r_i](r_i \times \omega_b) \\
&= -[\omega_b](r_i \times (r_i \times \omega_b)) \\
&= -[\omega_b]((r_i \cdot \omega_b) r_i - (r_i \cdot r_i) \omega_b) \\
&= -[\omega_b]((r_i \cdot \omega_b) r_i - \|r_i\|^2 \omega_b) \\
&= -(r_i \cdot \omega_b)[\omega_b] r_i + \|r_i\|^2 [\omega_b] \omega_b \\
&= -(r_i \cdot \omega_b)(\omega_b \times r_i) + 0 \\
&= -(r_i \cdot \omega_b)(\omega_b \times r_i)
\end{align}

\textbf{因此}：$[r_i][\omega_b]^2 r_i = -[\omega_b][r_i]^2 \omega_b$。

\textbf{方法2：使用$[r_i]^2$的表达式}

\textbf{关键}：$[r_i]^2 = r_i r_i^T - \|r_i\|^2 I$

\textbf{因此}：
\begin{align}
[r_i][\omega_b]^2 r_i &= [r_i](\omega_b \omega_b^T - \|\omega_b\|^2 I) r_i \\
&= [r_i](\omega_b \omega_b^T) r_i - \|\omega_b\|^2 [r_i] r_i \\
&= [r_i](\omega_b \omega_b^T) r_i \quad \text{（因为$[r_i] r_i = 0$）} \\
&= r_i \times (\omega_b (\omega_b^T r_i)) \\
&= (\omega_b^T r_i)(r_i \times \omega_b) \\
&= -(\omega_b^T r_i)(\omega_b \times r_i)
\end{align}

\textbf{另一方面}：
\begin{align}
-[\omega_b][r_i]^2 \omega_b &= -[\omega_b](r_i r_i^T - \|r_i\|^2 I) \omega_b \\
&= -[\omega_b](r_i r_i^T) \omega_b + \|r_i\|^2 [\omega_b] \omega_b \\
&= -[\omega_b](r_i r_i^T) \omega_b \\
&= -\omega_b \times (r_i (r_i^T \omega_b)) \\
&= -(r_i^T \omega_b)(\omega_b \times r_i) \\
&= -(\omega_b^T r_i)(\omega_b \times r_i)
\end{align}

\textbf{因此}：$[r_i][\omega_b]^2 r_i = -[\omega_b][r_i]^2 \omega_b$。

\section{总结}

\subsection{主要恒等式}

\begin{enumerate}
\item \textbf{恒等式1}：
\begin{equation}
[r_i][\dot{\omega}_b] r_i = -[r_i]^2 \dot{\omega}_b
\end{equation}

\item \textbf{恒等式2}：
\begin{equation}
[r_i][\omega_b]^2 r_i = -[\omega_b][r_i]^2 \omega_b
\end{equation}
\end{enumerate}

\subsection{应用}

这些恒等式用于推导单刚体动力学方程中的力矩公式：
\begin{equation}
m_b = \sum_i m_i [r_i] \dot{v}_b + I_b \dot{\omega}_b + [\omega_b] I_b \omega_b
\end{equation}

其中$I_b = -\sum_i m_i [r_i]^2$是惯性矩阵。

\end{document}

